\documentclass[11pt]{article}

\usepackage[margin=1in]{geometry}
\usepackage{graphicx}
\usepackage{float}
\usepackage{amsmath}
\usepackage{amssymb}
\usepackage{booktabs}
\usepackage{xcolor}
\usepackage{hyperref}
\usepackage{enumitem}
\usepackage{caption}

\hypersetup{colorlinks=true, linkcolor=blue, urlcolor=blue, citecolor=blue}

\title{\textbf{Progress Update}}
\date{\today}

\begin{document}
\maketitle

\section{Input}

\textbf{Depth recording.} RGB + depth topics present, no gripper topic,
wrench topic present. 1013 RGB frames, 1013 depth frames. Scene: plate,
screwdriver (clamped), charger, Rubik's cube.

\textbf{Recording A.} Gripper topic present, wrench topic present but 0
messages. 1745 frames, 116.8\,s. Scene: fastener/knob insertion into a
pegboard mount.

\textbf{Recording B.} Same sensor profile as A. 726 frames, 49.0\,s. Same
scene as A.

Standard exporter cannot decode the ZED compressed-depth transport; all
three recordings extracted with a custom \texttt{.mcap} decoder (no ROS\,2
install required, decodes pose/wrench/image messages directly from the
bag's embedded schemas, decodes the compressed-depth transport by hand:
12-byte header + PNG, standard quantization formula for true depth).

\section{Experiments}

\begin{enumerate}[noitemsep]
  \item Event detection on the depth recording: force-magnitude peaks
        grouped by time-gap into activity clusters (no gripper topic, so
        no grasp/release).
  \item Seed placement + SAM\,2 propagation on each object identified by
        the clusters above (plate, screwdriver, charger), independently.
  \item Two automated fixes attempted for the seed-picking failure on the
        plate (see Results): relative-size scoring, depth-based
        near-camera prior.
  \item Label-free object proposal (DINOv2 attention $\times$ depth) on
        the depth recording, using real sensor depth instead of monocular
        depth estimation, across all 5 activity-cluster frames.
  \item Cross-recording object-identity clustering: pooled appearance
        embeddings from the depth recording and the original recording
        into one clustering run.
  \item Event detection extended to handle recordings with a gripper
        topic but zero wrench messages (recordings A, B).
  \item Seed placement + bidirectional SAM\,2 propagation on the grasped
        object, recordings A and B.
  \item Bidirectional propagation (previously forward-only) applied
        retroactively to the grasped object in the first recording.
\end{enumerate}

\section{Results}

\textbf{Depth recording, event detection.} 16 force peaks, grouped into 5
activity clusters mapping to 4 distinct objects plus one compound action
(charger moved toward the screwdriver clamp).

\textbf{Depth recording, seed placement.} Automatic seed picker's area
cutoff (rejects any candidate mask $>$40\% of frame) rejected the correct
plate mask outright; picker fell back to a background sliver, which then
propagated as an unstable mask (17--78\% of frame) across all 1013 frames.
Manually reseeded; corrected propagation peaks at 31\% of frame at the
press and decays cleanly through occlusion.

\textbf{Automated seed-picking fixes.} Both failed. Relative-size scoring
fixes the depth recording but mis-picks a reflective table surface as the
target object on the first recording. Depth-based near-camera prior fails
because depth-sensor dropout is worst on the same specular surfaces
causing the original confusion. \textbf{Mitigation:} calibrated
end-effector projection (camera intrinsics + extrinsics) replaces the
heuristic entirely -- see Asks.

\textbf{N-object propagation.} Plate, screwdriver, and charger each
independently seeded and propagated across all 1013 frames, no runaway on
any (max 43,000\,px / 46,700\,px for the two newly added objects). All
three verified simultaneously masked in one frame with no cross-object
contamination.

\begin{figure}[H]
\centering
\includegraphics[width=0.7\textwidth]{figures/identify_depth_multi/overlays/f0786_1782835535683396106.png}
\caption{Plate (magenta), screwdriver (orange), charger (yellow) --
independently propagated, composed in one frame.}
\end{figure}

\textbf{Real depth vs.\ estimated depth, label-free proposal.} Depth-sensor
validity: 16\% on the specular plate, 58--72\% on matte tool/charger
surfaces. Object proposal correctly lands on the target object for 2 of 5
events; hits background scatter on the other 3 -- fails even with real
depth. \textbf{Mitigation:} none needed; confirms proprioceptive event
cueing is required and cannot be replaced by a label-free proposer, real
depth or not.

\textbf{Cross-recording object identity.} 3 true objects fragment into 12
clusters at default settings, 2 clusters contaminated with mixed objects.
Sweep across crop style / clustering threshold: keeping background context
in the crop (vs.\ tight, background-removed crops) eliminates
contamination across a threshold range, does not fix fragmentation.
\textbf{Mitigation:} contamination fixed (use background-context crops);
fragmentation unresolved, no fix attempted yet.

\textbf{Recordings A, B.} Grasp/release detected cleanly (A:
$t{=}39.47$\,s / $54.66$\,s; B: $t{=}10.69$\,s / $37.50$\,s), no press
event on either (0 wrench messages). Seed placement clean on both,
verified visually. Bidirectional propagation: full frame coverage on
both (1745/1745, 726/726), max mask area $<$1.5\% of frame, no drift.
Verified against a frame where the object is genuinely out of view (mask
correctly reads 0).

\begin{figure}[H]
\centering
\includegraphics[width=0.55\textwidth]{figures/t004/propagation_grasped/f0872_1782993012118220374.png}
\caption{Fastener tracked after insertion into the pegboard mount,
recording A.}
\end{figure}

\section{Asks}

\begin{enumerate}[noitemsep]
  \item Camera intrinsics (K, distortion) and camera-to-base extrinsics --
        unblocks calibrated end-effector projection, replaces the
        seed-picking heuristic.
  \item F/T sensor mount transform (sensor frame to base) -- needed for
        force-direction projection.
  \item Confirm F/T sensor status on recordings A and B (0 wrench
        messages on both -- no press/contact event detectable).
\end{enumerate}

\end{document}
